\begin{sgastatement}{\SGAInlineRefTarget{sga3:xxvi:corollary:6-9}Corollary 6.9}
Let \(S\) be a connected semilocal scheme, \(G\) an
\(S\)-reductive group, and \(Q\) a critical subtorus of \(G\), that
is, a subtorus such that
\[
 \rad\bigl(\Centr_G(Q)\bigr)=Q.
\]
The group \(\Centr_G(Q)\) is a Levi subgroup of a parabolic subgroup
of \(G\) if and only if
\[
 \Centr_G(Q)=\Centr_G(Q_{\mathrm{spl}}),
\]
or equivalently,
\[
 \operatorname{Lie}(G)^Q
   =\operatorname{Lie}(G)^{Q_{\mathrm{spl}}}.
\]
\end{sgastatement}
\SGARefTarget{sga3:xxvi:cor:6:9}

This follows from \SGARefLink{sga3:xxvi:cor:6:2}{6.2} and \SGARefLink{sga3:xxvi:prop:6:8}{6.8}.

\begin{sgastatement}{\SGAInlineRefTarget{sga3:xxvi:corollary:6-10}Corollary 6.10}
Let \(S\) be a connected semilocal scheme, \(G\) an
\(S\)-reductive group, and \(L\) a subgroup of \(G\).  The following
conditions are equivalent:
\begin{enumerate}[label=\textup{(\roman*)}]
\item There is a parabolic subgroup of \(G\) having \(L\) as a Levi
subgroup.
\item There is a split subtorus of \(G\) whose centralizer is \(L\).
\item There is a homomorphism
\[
 \mathbf G_{m,S}\longrightarrow G
\]
whose centralizer is \(L\).
\end{enumerate}
\end{sgastatement}
\SGARefTarget{sga3:xxvi:cor:6:10}

The implication \textup{(i)}\(\Rightarrow\)\textup{(ii)} follows
from \SGARefLink{sga3:xxvi:prop:6:8}{6.8}, and \textup{(iii)}\(\Rightarrow\)\textup{(i)} from \SGARefLink{sga3:xxvi:proposition:6-1}{6.1}.
It remains to prove
\textup{(ii)}\(\Rightarrow\)\textup{(iii)}.  Suppose therefore that
\[
 L=\Centr_G(Q),\qquad Q=D_S(M).
\]
Write \(\mathfrak g=\operatorname{Lie}(G/S)\) and
\[
 \mathfrak g=\bigoplus_{\alpha\in M}\mathfrak g^\alpha,
\]
and let \(R\) be the set of
\(\alpha\in M\setminus\{0\}\) for which
\(\mathfrak g^\alpha\ne0\).

Since \(R\) is finite and does not contain \(0\), there is a
homomorphism
\[
 u:M\longrightarrow\mathbf Z
\]
such that \(u(\alpha)\ne0\) for every \(\alpha\in R\).  By duality,
\(u\) yields a homomorphism \(\mathbf G_{m,S}\to Q\), and hence a
homomorphism
\[
 f:\mathbf G_{m,S}\longrightarrow G.
\]
One has
\[
 \Centr_G(f)\supset\Centr_G(Q).
\]
These are two smooth subgroups of \(G\), with connected fibers, and
their Lie algebras coincide, since both are equal to
\(\mathfrak g^0\).  They are therefore equal, by the usual
argument.

\begin{sgastatement}{\SGAInlineRefTarget{sga3:xxvi:corollary:6-11}Corollary 6.11}
Let \(S\) be a connected semilocal scheme and \(G\) an
\(S\)-reductive group.  The maps
\[
 L\longmapsto\rad(L)_{\mathrm{spl}},
 \qquad
 Q\longmapsto\Centr_G(Q)
\]
are mutually inverse, order-reversing bijections between:
\begin{itemize}
\item the set of subgroups \(L\) of \(G\) which are Levi subgroups
of parabolic subgroups of \(G\); and
\item the set of split subtori \(Q\) of \(G\) such that
\[
 \rad\bigl(\Centr_G(Q)\bigr)_{\mathrm{spl}}=Q.
\]
\end{itemize}
\end{sgastatement}
\SGARefTarget{sga3:xxvi:cor:6:11}

\begin{sgastatement}{Corollary 6.12}
Let \(S\) be a connected semilocal scheme and \(G\) an
\(S\)-reductive group.  Consider the following assertions:
\begin{enumerate}[label=\textup{(\roman*)}]
\item \(G\) has a parabolic subgroup different from \(G\).
\item \(G\) has a noncentral split subtorus.
\item[\textup{(ii bis)}] \(G\) has a noncentral split subtorus of
relative dimension \(1\).
\item There is a group homomorphism
\[
 \mathbf G_{a,S}\longrightarrow G
\]
which is a closed immersion.
\end{enumerate}
Then
\[
 \textup{(i)}
 \Longleftrightarrow\textup{(ii)}
 \Longleftrightarrow\textup{(ii bis)}
 \Longrightarrow\textup{(iii)}.
\]
\end{sgastatement}
\SGARefTarget{sga3:xxvi:cor:6:12}

The only new assertion is
\textup{(i)}\(\Rightarrow\)\textup{(iii)}.  Let \(P\ne G\) be a
parabolic subgroup of \(G\).  Then \(U=\rad^u(P)\ne e\).  Consider
the last nontrivial subgroup \(U_n\) in the composition series of
\(U\), as in \SGARefLink{sga3:xxvi:proposition:2-1}{2.1}.  There is an isomorphism
\[
 U_n\simeq W(E_n),
\]
where \(E_n\) is a locally free \(\mathcal O_S\)-module, and hence
is free because \(S\) is semilocal and connected.  Since
\(E_n\ne0\), there is a monomorphism
\[
 \mathcal O_S\longrightarrow E_n
\]
whose image is locally a direct summand.  It induces a closed
immersion
\[
 \mathbf G_{a,S}=W(\mathcal O_S)
   \lhook\joinrel\longrightarrow W(E_n)\simeq U_n,
\]
which proves \textup{(iii)}.

\begin{sgastatement}{Remark 6.12.1}
When \(S\) is the spectrum of a field of characteristic \(0\), the
Jacobson--Morozov theorem gives
\textup{(iii)}\(\Rightarrow\)\textup{(ii bis)}.  The four preceding
conditions are then equivalent \textup{(``Godement's criterion'';
cf.\ \textup{[BT65]}, 8.5)}.\footnote{More generally, this remains
true when \(S\) is the spectrum of a perfect field \textup{(Tits)};
see \textup{[BT71]}, Corollary 3.7.}
\end{sgastatement}
\SGARefTarget{sga3:xxvi:rem:6:12:1}

\begin{sgastatement}{Definition 6.13}
Let \(S\) be a connected semilocal scheme and \(G\) an
\(S\)-reductive group.  The group \(G\) is said to be
\emph{anisotropic} if it contains no nontrivial split subtorus.
\end{sgastatement}
\SGARefTarget{sga3:xxvi:def:6:13}

\begin{sgastatement}{Corollary 6.14}
Let \(S\) be a connected semilocal scheme.  An \(S\)-reductive group
\(G\) is anisotropic if and only if it has no parabolic subgroup
\(P\ne G\) and its radical is anisotropic.
\end{sgastatement}
\SGARefTarget{sga3:xxvi:cor:6:14}

Combining \SGARefLink{sga3:xxvi:lemma:6-6}{6.6} with \SGARefLink{sga3:xxiv:theorem:4-1-5}{Exposé XXIV, 4.1.5}, and \SGARefLink{sga3:xxii:subsection:6-2}{Exposé XXII, 6.2} gives:

\begin{sgastatement}{Corollary 6.15}
Let \(S\) be a connected semilocal scheme and \(G\) an isotrivial
\(S\)-reductive group \textup{(for example, \(G\) semisimple, or
\(S\) normal; cf.\ \SGARefLink{sga3:x:theorem:5-16}{Exposé X, 5.16})}.  The group \(G\) is anisotropic
if and only if it has no parabolic subgroup \(P\ne G\) and
\[
 \operatorname{Hom}_{S\text{-}\mathrm{gr}}
   (G,\mathbf G_{m,S})=e.
\]
\end{sgastatement}
\SGARefTarget{sga3:xxvi:cor:6:15}

\begin{sgastatement}{Proposition 6.16}
Let \(S\) be a connected semilocal scheme and \(G\) an
\(S\)-reductive group.  The maximal split subtori of \(G\) are
precisely the largest central split subtori of the Levi subgroups of
the minimal parabolic subgroups of \(G\).  Any two such tori are
conjugate by an element of \(G(S)\).
\end{sgastatement}
\SGARefTarget{sga3:xxvi:prop:6:16}

Let \(Q\) be a maximal split subtorus of \(G\).  By \SGARefLink{sga3:xxvi:cor:6:2}{6.2},
\[
 L=\Centr_G(Q)
\]
is a Levi subgroup of a parabolic subgroup \(P\) of \(G\).  Since
\[
 Q\subset\rad\bigl(\Centr_G(Q)\bigr)_{\mathrm{spl}},
\]
the maximality of \(Q\) gives
\[
 Q=\rad\bigl(\Centr_G(Q)\bigr)_{\mathrm{spl}}.
\]
By \SGARefLink{sga3:xxvi:corollary:6-11}{6.11}, \(L\) is minimal among the Levi subgroups of parabolic
subgroups of \(G\).  Hence \(P\) is a minimal parabolic subgroup of
\(G\), by \SGARefLink{sga3:xxvi:proposition:1-20}{1.20}.  It follows from \SGARefLink{sga3:xxvi:corollary:5-7}{5.7} and \SGARefLink{sga3:xxvi:corollary:5-5}{5.5}\textup{(iv)} that any
two such subtori \(Q\) are conjugate by an element of \(G(S)\).
The conjugacy of the tori \(Q\) and of the pairs \((P,L)\) also
proves the first assertion.

\begin{sgastatement}{Corollary 6.17}
Let \(S\) be a connected semilocal scheme, and let \(P,P'\) be two
minimal parabolic subgroups in standard position
\textup{(\SGARefLink{sga3:xxvi:def:4:5:1:1}{4.5.1.1})}.  Then \(P\cap P'\) contains a common Levi
subgroup of \(P\) and \(P'\).
\end{sgastatement}
\SGARefTarget{sga3:xxvi:cor:6:17}

Indeed, \(P\cap P'\) contains a maximal torus \(T\) of \(G\), by
\SGARefLink{sga3:xxvi:proposition:4-5-1}{4.5.1}\textup{(v)}.  Let \(L\) be the unique Levi subgroup of \(P\)
which contains \(T\).  By \SGARefLink{sga3:xxvi:proposition:1-21}{1.21},
\[
 \rad(P)\cap T=\rad(P)\cap L=\rad(L).
\]
Thus \(\rad(P)\cap T\) contains
\(\rad(L)_{\mathrm{spl}}\), which is a maximal split subtorus of
\(G\), and is therefore necessarily equal to
\(T_{\mathrm{spl}}\).  Hence
\[
 L=\Centr_G(T_{\mathrm{spl}}),
\]
and, by symmetry, \(L\) is also a Levi subgroup of \(P'\).

\begin{sgastatement}{Remark 6.18}
It follows from \SGARefLink{sga3:xxvi:proposition:1-21}{1.21} that a parabolic subgroup \(P\) of \(G\) is
minimal if and only if \(\rad(P)\) contains a maximal split subtorus
of \(G\).  In that case, if \(T\) is a maximal torus of \(P\), the
torus \(T_{\mathrm{spl}}\) is a maximal split subtorus both of \(G\)
and of \(\rad(P)\), and \(\Centr_G(T_{\mathrm{spl}})\) is a Levi
subgroup of \(P\).  Moreover, every Levi subgroup of \(P\) is
obtained in this way.
\end{sgastatement}
\SGARefTarget{sga3:xxvi:rem:6:18}
